\documentclass[../../main.tex]{subfiles}% 注意这里的文件路径不能用 ./main.tex ,否则用latexmk编译子文件会报错
\graphicspath{{\subfix{./image/}}{\subfix{../../image/}}} % 指定图片引用路径,后续可以直接使用图片文件名
% 如果图片存放在上级目录,则使用 ../../image/ 作为备用路径

% 例如:
% \begin{figure}[H]
% \centering
% \includegraphics[width=0.3\textwidth]{图.png}
% \caption{}
% \label{figure:图}
% \end{figure}
% 注意:上述\label{}一定要放在\caption{}之后,否则引用图片序号会只会显示??.

\begin{document}

\section{函数性态分析综合}

\begin{proposition}
设 $f : (a, b) \to (a, b)$ 满足对任意的 $x, y \in (a, b)$, 当 $x \neq y$ 时, 有 
$$
\left| f(x) - f(y) \right| < \left| x - y \right|.
$$ 
任取 $x_1 \in (a, b)$, 令 $$
x_{n+1} = f(x_n), n = 1, 2, \cdots.
$$
证明: 数列 $\{x_n\}_{n=1}^{\infty}$ 收敛.
\end{proposition}
\begin{proof}
注意到 $x_1 \in (a, b)$，假设 $x_k \in (a, b)$，则 $x_{k+1} = f(x_k) \in (a, b)$，故由数学归纳法可知 $x_n \in (a, b), \forall n \in \mathbb{N}$。
又由条件可知，对 $\forall \varepsilon > 0$，令 $\delta = \varepsilon > 0$，当 $x, y \in (a, b)$ 且 $0 < |x - y| < \delta$ 时，有
\begin{align*}
|f(x) - f(y)| < |x - y| < \delta = \varepsilon。
\end{align*}
故 $f$ 在 $(a, b)$ 上一致连续。从而 $f \in C(a, b)$.由Cauchy收敛准则知,$\underset{x\rightarrow a^+}{\lim}f\left( x \right) ,\underset{x\rightarrow b^-}{\lim}f\left( x \right) $存在,因此可将$f$连续延拓到$[a,b]$.令 $F(x) = f(x) - x$，则$F \in C[a, b]$.下面我们对 $F$ 进行分类讨论。

\begin{enumerate}[(1)]
\item 若 $F$ 在 $(a, b)$ 上不变号，则由 $F \in C(a, b)$ 可知，$F$ 要么恒大于零，要么恒小于零。不妨设 $F$ 在 $(a, b)$ 上恒大于零，即 $f(x) > x, \forall x \in (a, b)$。从而
\begin{align*}
x_{n+1} = f(x_n) > x_n, \forall n \in \mathbb{N}。
\end{align*}
即 $\{x_n\}$ 单调递增。又因为 $x_n \in (a, b), \forall n \in \mathbb{N}$，所以由单调有界定理可知 $\{x_n\}$ 收敛。

\item 若 $F$ 在 $(a, b)$ 上变号，则由 $F \in C(a, b)$ 及介值定理可得，存在 $\xi \in (a, b)$，使得 $f(\xi) = \xi$。
若存在 $\xi' \in (a, b)$ 且 $\xi' \neq \xi$，使得 $f(\xi') = \xi'$，则由条件可得到
\begin{align*}
|\xi - \xi'| = |f(\xi) - f(\xi')| < |\xi - \xi'|。
\end{align*}
显然矛盾！因此存在唯一的 $\xi \in (a, b)$，使得 $f(\xi) = \xi$。从而
\begin{align*}
|x_{n+1} - \xi| = |f(x_n) - f(\xi)| < |x_n - \xi|, \forall n \in \mathbb{N}。
\end{align*}
于是 $\{|x_n - \xi|\}$ 单调递减且有下界 $0$，故由单调有界定理可知 $\lim_{n \to \infty} |x_n - \xi| = A \geqslant 0$ 存在。
\begin{enumerate}[(i)]
\item 当 $A = 0$ 时，则由 $\lim_{n \to \infty} |x_n - \xi| = A = 0$ 可得 $\lim_{n \to \infty} x_n = \xi$。
\item 当 $A > 0$ 时，若 $\{x_n\}$ 收敛，则结论已经成立。若 $\{x_n\}$ 发散，则由 $x_n \in (a, b), \forall n \in \mathbb{N}$ 及聚点定理可知，$\{x_n\}$ 至少有一个聚点。若 $\{x_n\}$ 只有一个聚点，则 $\{x_n\}$ 收敛与假设矛盾！因此 $\{x_n\}$ 至少有两个聚点。任取收敛子列 $\{x_{n_k}\} \subset \{x_n\}$，设 $\lim_{k \to \infty} x_{n_k} = B$，则
\begin{align*}
A=\lim_{n\rightarrow \infty} |x_n-\xi |=\lim_{k\rightarrow \infty} |x_{n_k}-\xi |=|B-\xi |.
\end{align*}
从而 $B = \xi - A$ 或 $\xi + A$。因此 $\{x_n\}$ 最多有两个聚点 $\xi - A,\,\xi + A\in [a,b]$。故$\{x_n\}$ 有且仅有两个聚点 $\xi - A$ 和 $\xi + A$。进而一定存在收敛子列 $\{x_{n_k}\}$，使得$\lim_{k \to \infty} x_{n_k} = \xi - A.$因为$\xi-A\ne \xi$不是$f$的不动点,而$\{x_n\}$只有两个聚点,所以
\begin{align*}
\lim_{k\rightarrow \infty} x_{n_k+1}=\lim_{k\rightarrow \infty} f\left( x_{n_k} \right) =f\left( \xi -A \right) \ne \xi -A\Longrightarrow \lim_{k\rightarrow \infty} x_{n_k+1}=f(\xi -A)=\xi +A.
\end{align*}
由
\begin{align*}
\lim_{k\rightarrow \infty}x_{n_k}=\xi -A<\xi ,\quad \lim_{k\rightarrow \infty}x_{n_k+1}=\xi +A>\xi 
\end{align*}
知，存在$K\in \mathbb{N}$，使得$\forall k>K$，有$x_{n_k}<\xi$，$x_{n_k+1}>\xi$。又$\{| x_n-\xi |\}$递减趋于$A$，故对$\forall k>K$有
\begin{gather*}
A\leqslant | x_{n_k}-\xi |=\xi -x_{n_k}<\xi -a\Longrightarrow \xi -A>a,
\\
A\leqslant | x_{n_k+1}-\xi |=x_{n_k+1}-\xi <b-\xi \Longrightarrow \xi +A<b.
\end{gather*}
因此$\xi -A,\xi +A\in (a,b)$，$\xi =f(\xi)$。
再由条件可得
\begin{align*}
A = |\xi - (\xi + A)| = |f(\xi) - f(\xi - A)| < |\xi - (\xi - A)| = A。
\end{align*}
显然矛盾！故 $A > 0$ 不成立，于是 $A = 0$。再由 (1) 可得 $\lim_{n \to \infty} x_n = \xi$，即 $\{x_n\}$ 收敛，与假设$\{x_n\}$发散矛盾！
\end{enumerate}
\end{enumerate}
\end{proof}

\begin{proposition}[时滞方程]\label{proposition:时滞方程}
设 $f$ 在 $\mathbb{R}$ 上可微且满足
\begin{align*}
\lim_{x \to +\infty} f'(x) = 1, \quad
f(x+1) - f(x) = f'(x), \forall x \in \mathbb{R}.
\end{align*}
证明存在常数 $C \in \mathbb{R}$ 使得 $f(x) = x + C, \forall x \in \mathbb{R}$.
\end{proposition}
\begin{proof}
由 \( f(x+1)-f(x)=f'(x),\forall x\in \mathbb{R} \) 及 \( f\in D(\mathbb{R}) \) 可知 \( f' \in C(\mathbb{R}) \)。
对 \(\forall x_1\in \mathbb{R} \)，固定 \( x_1 \)，记
\[ A=\{ z>x_1 \mid f'(z) = f'(x_1) \} .\]
由 Lagrange 中值定理及 \( f(x+1)-f(x)=f'(x),\forall x\in \mathbb{R} \) 可知
\[ \exists x_2\in (x_1,x_1+1) \,\,\mathrm{s}.\mathrm{t}.\,\, f'(x_1) = f(x_1+1) - f(x_1) = f'(x_2) .\]
故 \( x_2 \in A \)，从而 \( A \) 非空。现在考虑 \( y \triangleq \mathrm{sup}A \in (x_1,+\infty) \)，下证 \( y=+\infty \)。
若 \( y<+\infty \)，则存在 \(\{ z_{n}' \}_{n=1}^{\infty} \)，使得
\[ z_{n}' \rightarrow y \text{且} f'(z_{n}') = f'(x_1) .\]
两边同时令 \( n \rightarrow \infty \)，由 \( f' \in C(\mathbb{R}) \) 可得
\[ f'(x_1) = \lim_{n \rightarrow \infty} f'(z_{n}') = f'(y) .\]
又由 Lagrange 中值定理及 \( f(x+1)-f(x)=f'(x),\forall x\in \mathbb{R} \) 可得
\[ \exists y' \in (y,y+1) \,\,\mathrm{s}.\mathrm{t}.\,\, f'(y) = f(y+1) - f(y) = f'(y') .\]
从而 \( y' \in A \) 且 \( y' > y \)，这与 \( y=\mathrm{sup}A \) 矛盾！故 \( y=+\infty \)。于是存在 \(\{ z_n \}_{n=1}^{\infty} \)，使得
\[ z_n \rightarrow +\infty \text{且} f'(z_n) = f'(x_1) .\]
两边同时令 \( n \rightarrow \infty \)，由 \( f' \in C(\mathbb{R}) \) 及 \(\lim_{x \rightarrow +\infty} f'(x) = 1 \) 可得
\[ f'(x_1) = \lim_{n \rightarrow \infty} f'(z_n) = \lim_{x \rightarrow +\infty} f'(x) = 1 .\]
因此由 \( x_1 \) 的任意性得，存在 \( C \) 为常数，使得 \( f(x) = x+C,\forall x\in \mathbb{R} \).
\end{proof}

\begin{example}
设 $f\in C^2(\mathbb{R})$ 满足 $f(1)\leqslant0$ 以及
\begin{align*}
\lim_{x\to\infty}[f(x)-|x|]=0.
\tag{12.27}
\end{align*}
证明:
(1)存在 $\xi\in(1,+\infty)$, 使得 $f'(\xi)>1$.

(2)存在 $\eta\in\mathbb{R}$, 使得 $f''(\eta)=0$.
\end{example}
\begin{proof}
(1)如果对任何 $x\in(1,+\infty)$, 都有 $f'(x)\leqslant1$, 那么 $[f(x)-x]'\leqslant0$ 知 $f(x)-x$ 在 $[1,+\infty)$ 单调递减. 从而
\begin{align*}
-1\geqslant f(1)-1\geqslant\lim_{x\to+\infty}[f(x)-x]=\lim_{x\to\infty}[f(x)-|x|]=0,
\end{align*}
这就是一个矛盾! 于是我们证明了 (1).

(2)
若对任何 $x\in\mathbb{R}$, 我们有 $f''(x)\neq0$.从而$f''(x)$要么恒大于零，要么恒小于零,否则由零点存在定理可得矛盾!任取$\xi \in \mathbb{R}$.

当 $f''(x)>0,\forall x\in\mathbb{R}$, 我们知道 $f$ 在 $\mathbb{R}$ 上是下凸函数. 由 (1) 和下凸函数切线总是在函数下方, 我们知道
\begin{align*}
f(x)\geqslant f(\xi)+f'(\xi)(x-\xi),\forall x>\xi.
\end{align*}
于是
\begin{align*}
0=\lim_{x\to+\infty}[f(x)-x]\geqslant\lim_{x\to+\infty}[f(\xi)-f'(\xi)\xi+(f'(\xi)-1)x]=+\infty,
\end{align*}
这就是一个矛盾!

当 $f''(x)<0,\forall x\in\mathbb{R}$, 我们知道 $f$ 在 $\mathbb{R}$ 上是上凸函数. 由 (1) 和上凸函数切线总是在函数上方, 我们有
\begin{align*}
f(x)\leqslant f(\xi)+f'(\xi)(x-\xi),\forall x<\xi.
\end{align*}
于是
\begin{align*}
0=\lim_{x\to-\infty}[f(x)+x]\leqslant\lim_{x\to-\infty}[f(\xi)-f'(\xi)\xi+(f'(\xi)+1)x]=-\infty,
\end{align*}
这就是一个矛盾! 因此我们证明了 (2). 
\end{proof}

\begin{example}
设 \(f\) 在 \([a,b]\) 上每一个点极限都存在，证明：\(f\) 在 \([a,b]\) 有界。
\end{example}
\begin{note}
极限存在必然局部有界，本题就是说局部有界可以推出在紧集上有界。在大量问题中会有一个公共现象：即\textbf{局部的性质等价于在所有紧集上的性质}。证明的想法就是有限覆盖。 
\end{note}
\begin{proof}
对 \(\forall c\in [a,b]\)，由 \(\lim_{x\rightarrow c}f(x)\) 存在可知，存在 \(c\) 的邻域 \(U_c\) 和 \(M>0\)，使得
\begin{align*}
\sup_{x\in U_c\cap [a,b]}|f(x)|\leqslant M_c.
\end{align*}
注意 \([a,b]\subset \bigcup_{c\in [a,b]}U_c\)，由有限覆盖定理得，存在 \(c_1,c_2,\cdots,c_n\in [a,b]\)，使得
\begin{align*}
[a,b]\subset \bigcup_{k = 1}^nU_{c_k}. 
\end{align*}
故 \(\sup_{x\in [a,b]}|f(x)|\leqslant \max_{1\leqslant k\leqslant n}M_k\)。 
\end{proof}

\begin{example}
设 \(f\) 是 \((a, +\infty)\) 有界连续函数, 证明对任何实数 \(T\) , 存在数列 \(\lim_{n \to \infty} x_n = +\infty\) 使得
\begin{align*}
\lim_{n \to \infty} [f(x_n + T) - f(x_n)] = 0.
\end{align*}
\end{example}
\begin{remark}
因为\(\vert f(x + T) - f(x)\vert \geqslant 0\) , 所以
\begin{align*}
0\leqslant \varliminf_{x \to +\infty} \vert f(x + T) - f(x)\vert \leqslant \varlimsup_{x \to +\infty} \vert f(x + T) - f(x)\vert.
\end{align*}
原结论的反面只用考虑\(\varliminf_{x \to +\infty} \vert f(x + T) - f(x)\vert\) 即可. 若\(\varliminf_{x \to +\infty} \vert f(x + T) - f(x)\vert = 0\) , 则一定存在子列 \(x_n \to +\infty\) , 使得结论成立.因此原结论等价于\(\varliminf_{x \to +\infty} \vert f(x + T) - f(x)\vert = 0\) .
故原结论的反面就是\(\varliminf_{x \to +\infty} \vert f(x + T) - f(x)\vert > 0\) .
\end{remark}
\begin{note}
考虑反证法之后,再进行定性分析(画$f(x)$的大致走势图),就容易找到矛盾.
\end{note}
\begin{proof}
当$T=0$时,显然存在这样的数列.不妨设$T>0,$
假设\(\varliminf_{x \to +\infty} \vert f(x + T) - f(x)\vert > 0\) , 则存在\(\varepsilon_0 > 0\) , \(X > 0\) , 使得
\begin{align}
\vert f(x + T) - f(x)\vert \geqslant \varepsilon_0>0,\quad \forall x \geqslant X \label{example0.4section05----1.1}
\end{align}
令\(g(x) \triangleq f(x + T) - f(x)\) , 则若存在 \(x_1, x_2 \geqslant X\) , 使得
\[g(x_1) = f(x_1 + T) - f(x_1) \geqslant \varepsilon_0 > 0 ,\quad g(x_2) = f(x_2 + T) - f(x_2) \leqslant -\varepsilon_0 < 0. \]
不妨设 \(x_1 < x_2\) , 由 \(g\) 连续及介值定理可知, 存在 \(x_3 \in (x_1, x_2)\) , 使得
\begin{align*}
g(x_3) = f(x_3 + T) - f(x_3) = 0
\end{align*}
与\eqref{example0.4section05----1.1} 式矛盾! 故 \(g(x) \triangleq f(x + T) - f(x)\) 在\([X, +\infty)\) 上要么恒大于\(\varepsilon_0\) , 要么恒小于\(\varepsilon_0\) . 于是不妨设
\begin{align}
f(x + T) - f(x) \geqslant \varepsilon_0,\quad \forall x \geqslant X .\label{example0.4section05----1.2}
\end{align}
从而对\(\forall k \in \mathbb{N}\) ,当 \(x \geqslant X\) 时, 有
\(x + (k - 1)T > X\) .
于是由\eqref{example0.4section05----1.2}式可得
\begin{align*}
f(x + kT) - f(x + (k - 1)T) \geqslant \varepsilon_0,\quad \forall x \geqslant X. 
\end{align*}
进而对上式求和可得, 对\(\forall n \in \mathbb{N}\) , 都有
\begin{align*}
\sum_{k = 1}^n [f(x + kT) - f(x + (k - 1)T)] = f(x + nT) - f(x) \geqslant n\varepsilon_0,\quad \forall x \geqslant X.
\end{align*}
任取 \(x_0 \geqslant X\) , 则
\begin{align*}
f(x_0 + nT) - f(x_0) \geqslant n\varepsilon_0,\quad \forall n \in \mathbb{N}.
\end{align*}
令 \(n \to \infty\) , 得\(\lim_{x \to +\infty} f(x) = +\infty\) . 这与 \(f\) 在\((a, +\infty)\) 上有界矛盾! 
\end{proof}

\begin{proposition}
\begin{enumerate}
\item 设 \(f_n\in C[a,b]\) 且关于 \([a,b]\) 一致的有
\begin{align*}
\lim_{n \to \infty}f_n(x)=f(x).
\end{align*}
证明: 对 \(\{x_n\}\subset [a,b]\), \(\lim_{n \to \infty}x_n = c\), 有
\begin{align*}
\lim_{n \to \infty}f_n(x_n)=f(c).
\end{align*}

\item 设 \(f_n(x):\mathbb{R}\to\mathbb{R}\) 满足对任何 \(x_0\in\mathbb{R}\) 和 \(\{x_n\}_{n = 1}^{\infty}\subset\mathbb{R}\), \(\lim_{n \to \infty}x_n = x_0\), 都有
\begin{align*}
\lim_{n \to \infty}f_n(x_n)=f(x_0),
\end{align*}
证明: \(f\in C(\mathbb{R})\). 
\end{enumerate}
\end{proposition}
\begin{proof}
\begin{enumerate}
\item 由\(f_n\)一致收敛到\(f(x)\)可知, 对\(\forall \varepsilon > 0\), 存在\(N_0\in \mathbb{N}\), 使得对\(\forall N\geqslant N_0\), 当\(n\geqslant N\)时, 对\(\forall x\in [a,b]\), 都有
\begin{align*}
|f_n(x) - f_N(x)| < \varepsilon.
\end{align*}
从而由上式可得
\begin{align*}
|f_n(x_n) - f(c)| &\leqslant |f_n(x_n) - f_N(x_n)| + |f_N(x_n) - f_N(c)| + |f_N(c) - f(c)|\\
&\leqslant \varepsilon + |f_N(x_n) - f_N(c)| + |f_N(c) - f(c)|.
\end{align*}
令\(n\rightarrow +\infty\), 由\(f\)的连续性及\(\lim_{n\rightarrow \infty}x_n = c\)可得
\begin{align*}
\varlimsup_{n\rightarrow \infty}|f_n(x_n) - f(c)| &\leqslant \varepsilon + |f_N(c) - f(c)|.
\end{align*}
再令\(N\rightarrow +\infty\), 由\(\lim_{n\rightarrow \infty}f_n(x) = f(x)\), \(\forall x\in [a,b]\)可知
\begin{align*}
\varlimsup_{n\rightarrow \infty}|f_n(x_n) - f(c)| &\leqslant \varepsilon.
\end{align*}
令\(\varepsilon \rightarrow 0^+\), 得\(\varlimsup_{n\rightarrow \infty}|f_n(x_n) - f(c)| \leqslant 0\). 故\(\lim_{n\rightarrow \infty}f_n(x_n) = f(c)\). 

\item 反证, 若\(f\)在\(x_0\in \mathbb{R}\)处不连续, 则存在\(\varepsilon_0 > 0\), 使得\(\forall m\in \mathbb{N}\), 存在\(y_m\in (x_0 - \frac{1}{m}, x_0 + \frac{1}{m})\), 使得
\begin{align}
|f(y_m) - f(x_0)| \geqslant \varepsilon_0. \label{equation--0.4section051.1}
\end{align}
令条件中的\(x_n= y_m\), \(\forall n\in \mathbb{N}\), 从而由条件可知\(\lim_{n\rightarrow \infty}|f_n(y_m) - f(y_m)| = 0\), \(m = 1,2,\cdots\),
故对\(\forall m\in \mathbb{N}\), 存在严格递增的数列\(n_m\rightarrow +\infty\), 使得
\begin{align}
|f_{n_m}(y_m) - f(y_m)| < \frac{\varepsilon_0}{2}. \label{equation--0.4section051.2}
\end{align}
从而由\eqref{equation--0.4section051.1}\eqref{equation--0.4section051.2}式可知, 对\(\forall m\in \mathbb{N}\), 都有
\begin{gather}
|f(y_{n_m}) - f(x_0)| \geqslant \varepsilon_0, \label{equation--0.4section051.3}\\
|f_{n_m}(y_{n_m}) - f(y_{n_m})| < \frac{\varepsilon_0}{2}. \label{equation--0.4section051.4}
\end{gather}
因此由\eqref{equation--0.4section051.3}\eqref{equation--0.4section051.4}式可得, 对\(\forall m\in \mathbb{N}\), 都有
\begin{align}
|f_{n_m}(y_{n_m}) - f(x_0)| &\geqslant |f(y_{n_m}) - f(x_0)| - |f_{n_m}(y_{n_m}) - f(y_{n_m})| \geqslant \varepsilon_0 - \frac{\varepsilon_0}{2} = \frac{\varepsilon_0}{2}. \label{equation--0.4section051.5}
\end{align}
注意到\(y_m\rightarrow x_0\), 于是\(y_{n_m}\rightarrow x_0\). 从而由已知条件可知\(\lim_{m\rightarrow \infty}f_{n_m}(y_{n_m}) = f(x_0)\),这与\eqref{equation--0.4section051.5}式矛盾! 故\(f\in C(\mathbb{R})\). 
\end{enumerate}
\end{proof}

\begin{example}
设 \(g\in C(\mathbb{R})\) 且以 \(T > 0\) 为周期, 且有
\begin{align}
f(f(x))=-x^3 + g(x).\label{equation0.5example-1.1}
\end{align}
证明:不存在 \(f\in C(\mathbb{R})\), 使得\eqref{equation0.5example-1.1}式成立. 
\end{example}
\begin{proof}
由\hyperref[proposition:连续的周期函数的基本性质]{连续的周期函数的基本性质}可知,存在$M>0,$使得$|g(x)|\leqslant M.$反证,假设存在 \(f\in C(\mathbb{R})\), 使得\eqref{equation0.5example-1.1}式成立. 则
\begin{align}
&\underset{x\rightarrow +\infty}{\lim}f\left( f\left( x \right) \right) =\underset{x\rightarrow +\infty}{\lim}\left( -x^3+g\left( x \right) \right) =-\infty ,\label{equation0.5example-2.1}
\\
&\underset{x\rightarrow -\infty}{\lim}f\left( f\left( x \right) \right) =\underset{x\rightarrow -\infty}{\lim}\left( -x^3+g\left( x \right) \right) =+\infty .\label{equation0.5example-2.2}
\end{align}
假设\(\lim_{x\rightarrow +\infty}f(x) = A\in \mathbb{R}\), 则存在\(x_n\rightarrow +\infty\), 使得\(f(x_n)\rightarrow A\). 从而由\eqref{equation0.5example-1.1}式可得
\begin{align*}
f(A) = \lim_{n\rightarrow \infty}f(f(x_n)) = \lim_{n\rightarrow \infty}(-x_{n}^{3}+g(x_n)) = -\infty.
\end{align*}
上式显然矛盾! 又因为\(f\in C(\mathbb{R})\), 所以\(\lim_{x\rightarrow +\infty}f(x) = +\infty\)或\(-\infty\).
否则, 当\(x\rightarrow +\infty\)时,\(f(x)\)振荡(上下极限不相等,取$K$为上下极限和的一半即可), 则由介值定理可知, 存在$K>0$,\(y_n\rightarrow +\infty\), 使得\(f(y_n) = K\), \(n = 1,2,\cdots\).
从而由\eqref{equation0.5example-2.1}式可知
\begin{align*}
-\infty = \lim_{x\rightarrow +\infty}f(f(x)) = \lim_{n\rightarrow \infty}f(f(y_n)) = f(K).
\end{align*}
显然矛盾!

(i)若\(\lim_{x\rightarrow +\infty}f(x) = +\infty\), 则
\begin{align*}
+\infty = \lim_{x\rightarrow +\infty}f(x) = f(+\infty) = \lim_{x\rightarrow +\infty}f(f(x)) = \lim_{x\rightarrow +\infty}[-x^3 + g(x)] = -\infty.
\end{align*}
显然矛盾!

(ii)若\(\lim_{x\rightarrow +\infty}f(x) = -\infty\), 则
\begin{align}
f(-\infty) = \lim_{x\rightarrow +\infty}f(f(x)) = \lim_{x\rightarrow +\infty}[-x^3 + g(x)] = -\infty. \label{equation0.5example-0.1}
\end{align}
从而对上式两边同时作用\(f\)可得
\begin{align}
f(-\infty) = f(f(-\infty)) = \lim_{x\rightarrow -\infty}[-x^3 + g(x)] = +\infty. \label{equation0.5example-0.2}
\end{align}
于是\eqref{equation0.5example-0.1}式与\eqref{equation0.5example-0.2}式显然矛盾!
综上,\(f\in C(\mathbb{R})\)的解不存在. 
\end{proof}

\begin{example}
\begin{enumerate}
\item 设 \(f\in C[0,+\infty)\) 是有界的。若对任何 \(r\in\mathbb{R}\)，都有 \(f(x)=r\) 在 \([0,+\infty)\) 只有有限个或者无根，证明:
$\lim_{x \to +\infty} f(x) \text{ 存在}.$

\item 设 \(f\in C(\mathbb{R})\)，\(n\) 是一个非 \(0\) 正偶数，使得对任何 \(y\in\mathbb{R}\)，都有 \(\{x\in\mathbb{R}:f(x)=y\}\) 是 \(n\) 元集。证明: 这样的 \(f\) 不存在。 
\end{enumerate}
\end{example}
\begin{proof}
\begin{enumerate}
\item 反证，设\(\lim\limits_{x \to +\infty} f(x)\)不存在，由\(f\)有界，可设\(\varlimsup_{x \to +\infty} f(x)=A > B=\varliminf_{x \to +\infty} f(x)\)。
任取\(C\in (B,A)\)，则由\(\varlimsup_{x \to +\infty} f(x)=A > C\)可知，存在\(x_1\geqslant 0\)，使得\(f(x_1)>C\)。
又由\(\varliminf_{x \to +\infty} f(x)=B < C\)可知，存在\(x_2 > x_1 + 1\)，使得\(f(x_2)<C\)。

于是再由\(\varlimsup_{x \to +\infty} f(x)=A > C\)可知，存在\(x_3 > x_2 + 1\)，使得\(f(x_3)>C\)。
又由\(\varliminf_{x \to +\infty} f(x)=B < C\)可知，存在\(x_4 > x_3 + 1\)，使得\(f(x_4)<C\)。
依此类推，可得递增数列\(\{x_n\}\)，使得
\begin{align*}
x_{n + 1} > x_n + 1, \quad f(x_{2n - 1}) > C, \quad f(x_{2n}) < C, \quad n = 1,2,\cdots.
\end{align*}
从而由\(f\in C[0,+\infty)\)及介值定理可得，对\(\forall n\in\mathbb{N}\)，存在\(y_n\in (x_{2n - 1},x_{2n})\)，使得\(f(y_n)=C\)，矛盾！ 

\item 设\(x_1<x_2<\cdots <x_n\)是\(f\)的所有零点，记\(x_0 = -\infty\)，\(x_{n + 1} = +\infty\)，
则由\(f\)的连续性及介值定理可知，\(f\)在\(( x_{i - 1},x_i )\)上不变号。
这里共有\(n + 1\)个区间，现在考虑\(( x_{i - 1},x_i )\)，\(i = 2,3,\cdots,n\)，这\(n - 1\)个区间。于是由抽屉原理可知，这\(n - 1\)个区间中必存在\(\frac{n}{2}\)区间，使\(f\)在这\(\frac{n}{2}\)个区间内都同号。

不妨设\(f\)在这\(\frac{n}{2}\)个区间内恒大于\(0\)，记\(f\)在\([ x_{i - 1},x_i ]\)上的最大值记为\(f( m_i ) \triangleq M_i>0\)，其中\(m_i\in ( x_{i - 1},x_i )\)，\(i = 2,3,\cdots,n\)。由介值定理知，至少存在\(c_i\in ( x_{i - 1},m_i )\)，\(c_{i}^{\prime}\in ( m_i,x_i )\)，使得
\begin{align*}
f( c_i ) =f( c_{i}^{\prime} ) =\frac{1}{2}\underset{2\leqslant  i\leqslant  n}{\min}M_i>0,i = 2,3,\cdots,n.
\end{align*}
注意到在\(( x_0,x_1 )\)，\(( x_n,x_{n + 1} )\)上\(f\)必不同号。否则，不妨设在\(( x_0,x_1 )\)，\(( x_n,x_{n + 1} )\)上\(f\)恒大于\(0\)，则由\(f\in C( \mathbb{R} )\)可知，存在\(M>0\)，使得
\(\left| f( x ) \right|<M,\forall x\in [ x_1,x_{n + 1} ].\)
从而\(f( x ) \geqslant -M,\forall x\in \mathbb{R} \)。这与对\(\forall y\in \mathbb{R} \)，\(f( x ) =y\)都有根矛盾！

不妨设\(f\)在\(( x_0,x_1 )\)上恒小于\(0\)，在\(( x_n,x_{n + 1} )\)上恒大于\(0\)，则\(f\)在\(( x_n,x_{n + 1} )\)上无上界。否则，存在\(K>\underset{2\leqslant  i\leqslant  n}{\max}M_i>0\)，使得\(f( x ) <K,\forall x\in ( x_n,x_{n + 1} )\)。又因为
\begin{align*}
f( x ) <0<K,\forall x\in ( x_0,x_1 ),\quad f( x ) \leqslant \underset{2\leqslant  i\leqslant  n}{\max}M_i<K,\forall x\in ( x_1,x_n ).
\end{align*}
所以\(f( x ) <K,\forall x\in \mathbb{R} \)。这与对\(\forall y\in \mathbb{R} \)，\(f( x ) =y\)都有根矛盾！

又\(f( x_n ) =0\)，故至少存在一个\(c\in ( x_n,x_{n + 1} )\)，使得\(f( c ) =\frac{1}{2}\underset{2\leqslant  i\leqslant  n}{\min}M_i>0\)。
综上，至少有\(n + 1\)个点使得\(f( x ) =\frac{1}{2}\underset{2\leqslant  i\leqslant  n}{\min}M_i>0\)。这与\(\{ x\in \mathbb{R} :f( x ) =\frac{1}{2}\underset{2\leqslant  i\leqslant  n}{\min}M_i \}\)是\(n\)元集矛盾！
\end{enumerate}
\end{proof}

\begin{example}
设 \(a,b > 1\) 且 \(f: \mathbb{R} \to \mathbb{R}\) 在 \(x = 0\) 邻域有界。若
\begin{align*}
f(ax) = bf(x),\quad \forall x \in \mathbb{R}
\end{align*}
证明：\(f\) 在 \(x = 0\) 连续。 
\end{example}
\begin{proof}
注意到
\begin{align*}
f(0) = bf(0) \Rightarrow f(0) = 0.
\end{align*}
由条件可得
\begin{align}
f(ax) = bf(x) \Rightarrow f(x) = \frac{f(ax)}{b} = \frac{f(a^2x)}{b^2} = \cdots = \frac{f(a^nx)}{b^n}, \forall n \in \mathbb{N}. \label{eq:100.1}
\end{align}
因为\(f\)在\(x=0\)邻域有界，所以存在\(\delta > 0\)，使得
\begin{align}
|f(x)| \leqslant M, \forall x \in (-\delta, \delta). \label{eq:100.2}
\end{align}
从而对\(\forall \varepsilon > 0\)，取\(N \in \mathbb{N}\)，使得
\begin{align}
\frac{M}{b^N} < \varepsilon. \label{eq:100.3}
\end{align}
于是当\(x \in \left( -\frac{\delta}{a^N}, \frac{\delta}{a^N} \right)\)时，结合\eqref{eq:100.1}\eqref{eq:100.2}\eqref{eq:100.3}式，我们有
\begin{align*}
|f(x)| = \left| \frac{f(a^Nx)}{b^N} \right| \leqslant \frac{M}{b^N} < \varepsilon.
\end{align*}
故\(\lim_{x \to 0} f(x) = f(0) = 0.\)
\end{proof}

\begin{example}
设 \(f \in C(\mathbb{R})\) 满足 \(f(x), f(x^2)\) 都是周期函数，证明：\(f\) 为常值函数. 
\end{example}
\begin{proof}
由\hyperref[proposition:连续周期函数必一致连续]{连续周期函数必一致连续}可知，\(f(x),f(x^2)\) 在 \(\mathbb{R}\) 上一致连续。
于是对任意满足 \(\lim_{n \to \infty} (x_n' - x_n'') = 0\) 的数列 \(\{x_n'\}, \{x_n''\}\)，都有
\begin{align}
|f(x_n') - f(x_n'')|, |f((x_n')^2) - f((x_n'')^2)| \to 0, \quad n \to \infty. \label{eq:101.1}
\end{align}
设 \(f(x)\) 的周期为 \(T > 0\)，则对 \(\forall c \in \mathbb{R}\)，取 \(x_n' = \sqrt{nT + c}, x_n'' = \sqrt{nT}\)，显然 \(x_n' - x_n'' = \frac{c}{\sqrt{nT + c} + \sqrt{nT}} \to 0\)。从而由 \eqref{eq:101.1} 式可得
\begin{align*}
|f((x_n')^2) - f((x_n'')^2)| = f(nT + c) - f(nT) = f(c) - f(0) \to 0.
\end{align*}
故 \(f(c) = f(0)\)，\(\forall c \in \mathbb{R}\)。故 \(f\) 为常值函数。
\end{proof}


\begin{example}
计算函数方程 \( f(\log_2 x) = f(\log_3 x) + \log_5 x \) 所有 \(\mathbb{R}\) 上的连续解.
\end{example}
\begin{note}
注意到
\[
f\left( \frac{\ln x}{\ln 2} \right) = f\left( \frac{\ln x}{\ln 3} \right) + \frac{\ln x}{\ln 5}, \, x > 0.
\]
为了凑裂项的形式, 我们待定
\[
f\left( \frac{\ln a_n}{\ln 2} \right) = f\left( \frac{\ln a_n}{\ln 3} \right) + \frac{\ln a_n}{\ln 5}, \, n \in \mathbb{N}.
\]
注意到我们有两种选择
\[
\frac{\ln a_n}{\ln 2} = \frac{\ln a_{n+1}}{\ln 3}, \quad \frac{\ln a_n}{\ln 3} = \frac{\ln a_{n+1}}{\ln 2}.
\]
前者公比 \(\frac{\ln 3}{\ln 2} > 1\), 后者公比 \(\frac{\ln 2}{\ln 3} < 1\), 为了求和方便我们选取后者.
\end{note}
\begin{proof}
设$f \in C(\mathbb{R})$，对$\forall x \in \mathbb{R}$，取$a_1 = x$，$\ln a_n = \left( \frac{\ln 2}{\ln 3} \right)^{n-1} \cdot \ln x$，$n \in \mathbb{N}$。则$\lim_{n \to \infty} \ln a_n = 0$。此时有
\begin{align*}
\frac{\ln a_n}{\ln 3} = \frac{\ln a_{n+1}}{\ln 2}, \forall n \in \mathbb{N}.
\end{align*}
于是由条件可得
\begin{align*}
f\left( \frac{\ln x}{\ln 2} \right) = f\left( \frac{\ln x}{\ln 3} \right) + \frac{\ln x}{\ln 5} \Rightarrow f\left( \frac{\ln a_n}{\ln 2} \right) = f\left( \frac{\ln a_n}{\ln 3} \right) + \frac{\ln a_n}{\ln 5}
\end{align*}
\begin{align*}
\Rightarrow f\left( \frac{\ln a_n}{\ln 2} \right) = f\left( \frac{\ln a_{n+1}}{\ln 2} \right) + \frac{\ln a_n}{\ln 5}, n = 1, 2, \cdots
\end{align*}
因此
\begin{align*}
\sum_{n=1}^{\infty} \left[ f\left( \frac{\ln a_n}{\ln 2} \right) - f\left( \frac{\ln a_{n+1}}{\ln 2} \right) \right] = \sum_{n=1}^{\infty} \frac{\ln a_n}{\ln 5} = \frac{1}{\ln 5} \cdot \frac{\ln x}{1 - \frac{\ln 2}{\ln 3}}.
\end{align*}
\begin{align*}
\sum_{n=1}^{\infty} \left[ f\left( \frac{\ln a_n}{\ln 2} \right) - f\left( \frac{\ln a_{n+1}}{\ln 2} \right) \right] = f\left( \frac{\ln a_1}{\ln 2} \right) - \lim_{n \to \infty} f\left( \frac{\ln a_{n+1}}{\ln 2} \right) = f\left( \frac{\ln x}{\ln 2} \right) - f(0).
\end{align*}
故
\begin{align*}
\frac{1}{\ln 5} \cdot \frac{\ln x}{1 - \frac{\ln 2}{\ln 3}} = f\left( \frac{\ln x}{\ln 2} \right) - f(0) \stackrel{y = \frac{\ln x}{\ln 2}}{\Rightarrow} f(y) = f(0) + \frac{y \ln 2 \ln 3}{\ln 5 \ln \frac{3}{2}}.
\end{align*}
\end{proof}

\begin{example}
设 \( n \in \mathbb{N} \)，\( f \in C^{n+2}(\mathbb{R}) \) 使得存在 \( \theta \in \mathbb{R} \) 满足对任何 \( x, h \in \mathbb{R} \) 都有
\begin{align*}
f(x + h) = f(x) + f'(x)h + \frac{f''(x)}{2}h^2 + \cdots + \frac{f^{(n-1)}(x)}{(n-1)!}h^{n-1} + \frac{f^{(n)}(x + \theta h)}{n!}h^n
\end{align*}
证明: \( f \) 是不超过 \( n + 1 \) 次的多项式.
\end{example}
\begin{proof}
对$\forall x,h\in \mathbb{R}$,由Taylor公式可知
\begin{align*}
f^{(n)}(x+\theta h) =f^{(n)}(x) +f^{(n+1)}(x) \theta h+\frac{f^{(n+1)}(\xi)}{2}\theta^2 h^2.
\end{align*}
再结合条件可得
\begin{align}
f(x+h) &=\sum_{j=0}^{n-1}\frac{f^{(j)}(x)}{j!}h^j+\frac{f^{(n)}(x+\theta h)}{n!}h^n \notag \\
&=\sum_{j=0}^n\frac{f^{(j)}(x)}{j!}h^j+\frac{f^{(n+1)}(x) \theta h+\frac{f^{(n+1)}(\xi)}{2}\theta^2 h^2}{n!}h^n, \label{eq:101.313}
\end{align}
由Taylor公式可知
\begin{align}
f(x+h) =\sum_{j=0}^{n+1}\frac{f^{(j)}(x)}{j!}h^j+\frac{f^{(n+2)}(\eta)}{(n+2)!}h^{n+2}. \label{eq:101.1313}
\end{align}
比较\eqref{eq:101.313}式和\eqref{eq:101.1313}式得
\begin{align}
\left[ \frac{1}{(n+1)!}-\frac{\theta}{n!} \right] f^{(n+1)}(x) =\left[ \frac{\theta^2 f^{(n+2)}(\xi)}{2n!}-\frac{f^{(n+2)}(\eta)}{(n+2)!} \right] h. \label{eq:101.12131}
\end{align}
当$\theta =\frac{1}{n+1}$时,我们有
\begin{align*}
\frac{\theta^2 f^{(n+2)}(\xi)}{2n!}=\frac{f^{(n+2)}(\eta)}{(n+2)!}.
\end{align*}
对上式令$h\rightarrow 0$,则$\xi,\eta \rightarrow x$,故此时就有
\begin{align*}
\frac{f^{(n+2)}(x)}{2n!(n+1)^2}=\frac{f^{(n+2)}(x)}{(n+2)!}\Rightarrow f^{(n+2)}(x) =0.
\end{align*}
当$\theta \ne \frac{1}{n+1}$时,对\eqref{eq:101.12131}式令$h\rightarrow 0$,则$\xi,\eta \rightarrow x$,故此时就有
\begin{align*}
\left[ \frac{1}{(n+1)!}-\frac{\theta}{n!} \right] f^{(n+1)}(x) =0\Rightarrow f^{(n+1)}(x) =0.
\end{align*}
因此,无论如何都有$f$是不超过$n+1$次的多项式.
\end{proof}

\begin{example}
\begin{enumerate}
\item 设
\begin{align}
P_n(x)=\frac{1}{2^n n!}\frac{\mathrm{d}^n}{\mathrm{d}x^n}(x^2 - 1)^n, \quad n = 1, 2, \cdots
\end{align}
证明多项式 \( P_n \) 的全部根都是实数且分布在 \( (-1, 1) \)。

\item 设 \( g(x)=e^{x^2}\frac{\mathrm{d}^n}{\mathrm{d}x^n}(e^{-x^2}) \)，证明 \( g \) 是只有实根的多项式。
\end{enumerate}
\end{example}
\begin{note}
本题第1问叫做Legendre(勒让德)多项式，第2问叫做Hermite多项式。第2问用Rolle定理时注意无穷远点也会提供零点。
\end{note}
\begin{proof}
\begin{enumerate}
\item 显然$P_n$是$n$次多项式,且$\pm 1$是$\frac{\mathrm{d}^k}{\mathrm{d}x^k}(x^2-1)^n$的$n-k$重根$(0\leqslant k\leqslant n)$.由Rolle定理可知,$\frac{\mathrm{d}}{\mathrm{d}x}(x^2-1)^n$在$(-1,1)$有一个实根.于是再由Rolle定理可知,$\frac{\mathrm{d}^2}{\mathrm{d}x^2}(x^2-1)^n$在$(-1,1)$有两个不同实根.反复利用Rolle定理可得,$\frac{\mathrm{d}^n}{\mathrm{d}x^n}(x^2-1)^n$在$(-1,1)$有$n$个不同实根.而$n$次多项式有且仅有$n$个根,故$P_n$的全部根都是实数且分布在$(-1,1)$上.

\item 设$\frac{\mathrm{d}^k}{\mathrm{d}x^k}(e^{-x^2}) = P_k(x)e^{-x^2}$，$P_k$是$k$次多项式，$k\in \mathbb{N}$，显然$P_0(x) = 1$，于是
\begin{align*}
\frac{\mathrm{d}^{k+1}}{\mathrm{d}x^{k+1}}(e^{-x^2}) = \left[ P_k'(x) - 2xP_k(x) \right] e^{-x^2}.
\end{align*}
令$P_{k+1}(x) = P_k'(x) - 2xP_k(x)$，则由$P_k$是$k$次多项式可知$P_{k+1}(x)$是$k+1$次多项式。
故由数学归纳法可知
\begin{align*}
\frac{\mathrm{d}^n}{\mathrm{d}x^n}(e^{-x^2}) = P_n(x)e^{-x^2}, \quad P_n \in \mathbb{R}[x], \quad n \in \mathbb{N}.
\end{align*}
因此$g(x) = e^{-x^2}\frac{\mathrm{d}^n}{\mathrm{d}x^n}(e^{-x^2}) = P_n(x)$是$n$次多项式$(n \in \mathbb{N})$。

显然$P_1(x)=-2x$只有一个实根$x=0$.
设$P_k$是有$k$个不同实根的多项式，这$k$个根为
\begin{align*}
x_1 < x_2 < \cdots < x_k.
\end{align*}
从而这$k$个根也是$P_k(x)e^{-x^2}$的根。由Rolle定理可知
\begin{align*}
P_{k+1}(x) = e^{-x^2}\frac{\mathrm{d}^k}{\mathrm{d}x^k}(e^{-x^2}) = e^{-x^2}\frac{\mathrm{d}}{\mathrm{d}x}\left( P_k(x)e^{-x^2} \right)
\end{align*}
在$(x_{j-1}, x_j)$，$j=2,3,\cdots,k$有实根。而$\lim\limits_{x \rightarrow \pm \infty} P_k(x)e^{-x^2} = 0$，故由\hyperref[theorem:加强的Rolle中值定理]{加强的Rolle中值定理}可知$P_{k+1}(x)$在$(-\infty, x_1)$，$(x_k, +\infty)$上还各有一个实根。因此$P_{k+1}(x)$有$k+1$个根。故由数学归纳法可知$P_n(x)$有$n$个实根$(n \in \mathbb{N})$。又因为$g(x) = P_n(x)$是$n$次多项式，而$n$次多项式有且仅有$n$个根，所以$g(x) = P_n(x)$是只有实根的多项式。
\end{enumerate}
\end{proof}

\begin{example}
设 $f$ 是直线上的非常值连续周期函数. 若 $g \in C(\mathbb{R})$ 且 $\varlimsup_{x \to +\infty} \frac{|g(x)|}{x} = +\infty$, 证明: $f \circ g$ 不是周期函数.
\end{example}
\begin{note}
$\varlimsup_{x \to +\infty} |g(x + \delta) - g(x)| = +\infty$.的证明类似\hyperref[theorem:函数Stolz定理]{函数Stolz定理}的证明.实际上就是利用了上极限版的函数Stolz定理,只不过我们之前并没有写出这个定理.
\end{note}
\begin{proof}
若 $f \circ g$ 是周期函数，则由\hyperref[proposition:连续周期函数必一致连续]{连续周期函数必一致连续}可知 $f \circ g$ 在 $\mathbb{R}$ 上一致连续.
设 $f$ 的周期为 $T > 0$, 记 $a \triangleq \max f - \min f > 0$, 则存在 $\delta > 0$, 使
\begin{align}
|f(g(x)) - f(g(y))| < a, \forall |x - y| \leqslant \delta. \label{eq:102.12}
\end{align}
先证 $\varlimsup_{x \to +\infty} |g(x + \delta) - g(x)| = +\infty$. 若 $\varlimsup_{x \to +\infty} |g(x + \delta) - g(x)| \ne +\infty$, 则存在 $A > 0$, 使
$|g(x + \delta) - g(x)| \leqslant A, \forall x \geqslant 0.$
对 $x \in [n\delta, (n + 1)\delta), n \in \mathbb{N}$, 我们有
\begin{align*}
\left| \frac{g(x)}{x} \right| \leqslant \frac{|g(x - n\delta)|}{n\delta} + \sum_{k=0}^{n - 1} \left| \frac{g(x - k\delta) - g(x - (k + 1)\delta)}{n\delta} \right| \leqslant \frac{1}{n\delta} \sup_{x \in [0, \delta]} |g(x)| + \frac{A}{\delta}.
\end{align*}
故 $\varlimsup_{x \to +\infty} \left| \frac{g(x)}{x} \right| \leqslant \frac{A}{\delta}$ 矛盾! 因此 $\varlimsup_{x \to +\infty} |g(x + \delta) - g(x)| = +\infty$.
于是存在 $x_0 \in \mathbb{R}$, 使得 $|g(x_0 + \delta) - g(x_0)| \geqslant T$. 由介值定理可知, 存在 $s, t \in [x_0, x_0 + \delta]$, 使得
$f(g(s)) = \max f, \quad f(g(t)) = \min f.$
从而由 \eqref{eq:102.12} 式可知
\begin{align*}
a = |f(g(s)) - f(g(t))| < a
\end{align*}
矛盾! 故 $f \circ g$ 不是周期函数 ($f \circ g$ 甚至不是一致连续函数).
\end{proof}

\begin{example}
设$F(x)$是$[0,+\infty)$上的单调递减函数，且
\begin{align}
\lim\limits_{x \to +\infty} F(x) = 0,\quad
\lim\limits_{n \to +\infty} \int_0^{+\infty} F(t) \sin \frac{t}{n} \mathrm{d}t = 0.\label{eq::23945789023589028225}
\end{align}
证明:
\begin{enumerate}[(i)]
\item \begin{align}
\lim\limits_{x \to +\infty} xF(x) = 0;\label{eq::23945789023589028226}
\end{align}

\item \begin{align}
\lim\limits_{x \to 0} \int_0^{+\infty} F(t) \sin(xt) \mathrm{d}t = 0.\label{eq::23945789023589028227}
\end{align}
\end{enumerate}
\end{example}
\begin{proof}
\begin{enumerate}[(i)]
\item 由\eqref{eq::23945789023589028225}知$F$非负. 由A-D判别法知本题涉及的积分都收敛. 注意到
\begin{align*}
\int_0^\infty F(t)\sin(tx)\mathrm{d}t &= \frac{1}{x}\int_0^\infty F\left(\frac{u}{x}\right)\sin u\mathrm{d}u = \frac{1}{x}\sum_{k=0}^\infty \int_{2k\pi}^{(2k+2)\pi} F\left(\frac{u}{x}\right)\sin u\mathrm{d}u \\
&= \frac{1}{x}\sum_{k=0}^\infty \left[ \int_{2k\pi}^{(2k+1)\pi} F\left(\frac{u}{x}\right)\sin u\mathrm{d}u + \int_{(2k+1)\pi}^{(2k+2)\pi} F\left(\frac{u}{x}\right)\sin u\mathrm{d}u \right] \\
&= \frac{1}{x}\sum_{k=0}^\infty \left[ \int_{2k\pi}^{(2k+1)\pi} F\left(\frac{u}{x}\right)\sin u\mathrm{d}u - \int_{2k\pi}^{(2k+1)\pi} F\left(\frac{u + \pi}{x}\right)\sin u\mathrm{d}u \right] \\
&= \frac{1}{x}\sum_{k=0}^\infty \int_{2k\pi}^{(2k+1)\pi} \left[ F\left(\frac{u}{x}\right) - F\left(\frac{u + \pi}{x}\right) \right]\sin u\mathrm{d}u \\
&\geqslant \frac{1}{x}\int_0^\pi \left[ F\left(\frac{u}{x}\right) - F\left(\frac{u + \pi}{x}\right) \right]\sin u\mathrm{d}u = \frac{1}{x}\int_0^{2\pi} F\left(\frac{u}{x}\right)\sin u\mathrm{d}u \geqslant 0,
\end{align*}
以及
\begin{align*}
\int_0^\infty F(t)\sin(tx)\mathrm{d}t &= \frac{1}{x}\int_0^\infty F\left(\frac{u}{x}\right)\sin u\mathrm{d}u = \frac{1}{x}\int_0^\pi F\left(\frac{u}{x}\right)\sin u\mathrm{d}u + \frac{1}{x}\sum_{k=0}^\infty \int_{(2k+1)\pi}^{(2k+3)\pi} F\left(\frac{u}{x}\right)\sin u\mathrm{d}u \\
&= \frac{1}{x}\int_0^\pi F\left(\frac{u}{x}\right)\sin u\mathrm{d}u + \frac{1}{x}\sum_{k=0}^\infty \left[ \int_{(2k+1)\pi}^{(2k+2)\pi} F\left(\frac{u}{x}\right)\sin u\mathrm{d}u + \int_{(2k+2)\pi}^{(2k+3)\pi} F\left(\frac{u}{x}\right)\sin u\mathrm{d}u \right] \\
&= \frac{1}{x}\int_0^\pi F\left(\frac{u}{x}\right)\sin u\mathrm{d}u + \frac{1}{x}\sum_{k=0}^\infty \left[ \int_{(2k+1)\pi}^{(2k+2)\pi} F\left(\frac{u}{x}\right)\sin u\mathrm{d}u - \int_{(2k+1)\pi}^{(2k+2)\pi} F\left(\frac{u + \pi}{x}\right)\sin u\mathrm{d}u \right] \\
&= \frac{1}{x}\int_0^\pi F\left(\frac{u}{x}\right)\sin u\mathrm{d}u + \frac{1}{x}\sum_{k=0}^\infty \int_{(2k+1)\pi}^{(2k+2)\pi} \left[ F\left(\frac{u}{x}\right) - F\left(\frac{u + \pi}{x}\right) \right]\sin u\mathrm{d}u \\
&\leqslant \frac{1}{x}\int_0^\pi F\left(\frac{u}{x}\right)\sin u\mathrm{d}u.
\end{align*}
从而
\begin{align}
0 \leqslant \int_0^{2\pi} \frac{F\left(\frac{t}{x}\right)\sin t}{x}\mathrm{d}t \leqslant \int_0^\infty F(t)\sin(tx)\mathrm{d}t = \int_0^\infty \frac{F\left(\frac{t}{x}\right)\sin t}{x}\mathrm{d}t \leqslant \int_0^\pi \frac{F\left(\frac{t}{x}\right)\sin t}{x}\mathrm{d}t. \label{eq::23945789023589028232}
\end{align}
上式取$x = \frac{1}{n}$并结合
\begin{align*}
\int_0^{+\infty} F(t) \sin \frac{t}{n} \mathrm{d}t &\geqslant n\int_0^{2\pi} F(nu)\sin u\mathrm{d}u \xlongequal{\text{区间再现}} n\int_0^\pi [F(nu) - F(2\pi n - nu)]\sin u\mathrm{d}u \\
&\geqslant n\int_0^{\frac{\pi}{2}} [F(nu) - F(2\pi n - nu)]\sin u\mathrm{d}u \geqslant n\int_0^{\frac{\pi}{2}} \left[ F\left(\frac{\pi n}{2}\right) - F\left(\frac{3\pi n}{2}\right) \right]\sin u\mathrm{d}u \\
&= n\left[ F\left(\frac{\pi n}{2}\right) - F\left(\frac{3\pi n}{2}\right) \right] \geqslant 0,
\end{align*}
和\eqref{eq::23945789023589028225}知
\[
\lim_{n \to +\infty} n\left[ F\left(\frac{\pi n}{2}\right) - F\left(\frac{3\pi n}{2}\right) \right] = 0.
\]
现在对任何$\varepsilon > 0$, 存在$N \in \mathbb{N}$使得对任何$n \geqslant N$都有
\begin{align}
n\left[ F\left(\frac{\pi n}{2}\right) - F\left(\frac{3\pi n}{2}\right) \right] \leqslant \varepsilon. \label{eq::23945789023589028233}
\end{align}
当正整数$k$充分大, 我们考虑
\[
b_n \triangleq k3^{n-1} \left[ F\left(\frac{\pi k3^{n-1}}{2}\right) - F\left(\frac{\pi k3^n}{2}\right) \right] \leqslant \varepsilon,
\]
则利用$\lim\limits_{x \to +\infty} F(x) = 0$和\eqref{eq::23945789023589028233}得
\[
0 \leqslant kF\left(\frac{k\pi}{2}\right) = k\sum_{n=1}^\infty \left[ F\left(\frac{\pi k3^{n-1}}{2}\right) - F\left(\frac{\pi k3^n}{2}\right) \right] = \sum_{n=1}^\infty \frac{b_n}{3^{n-1}} \leqslant \varepsilon \sum_{n=1}^\infty \frac{1}{3^{n-1}} = \frac{3}{2}\varepsilon.
\]
现在我们有$\lim\limits_{k \to +\infty} kF\left(\frac{k\pi}{2}\right) = 0$. 对$x \in [0,+\infty)$, 存在唯一的$k \in \mathbb{N}$使得$x \in \left[ \frac{k\pi}{2}, \frac{k+1}{2}\pi \right)$, 于是
\[
0 \leqslant xF(x) \leqslant \frac{(k+1)\pi}{2}F\left(\frac{k\pi}{2}\right) \to 0, k \to +\infty,
\]
因此我们证明了\eqref{eq::23945789023589028226}.

\item 我们由\eqref{eq::23945789023589028232}知
\begin{align}
\int_0^\infty F(t)\sin(tx)\mathrm{d}t \geqslant 0 \Rightarrow \varliminf_{x \to 0^+} \int_0^\infty F(t)\sin(tx)\mathrm{d}t \geqslant 0, \label{eq::23945789023589028234}
\end{align}
以及对任何$\eta > 0$, 我们有
\begin{align*}
\varlimsup_{x \to 0^+} \int_0^\infty F(t)\sin(tx)\mathrm{d}t &\leqslant \varlimsup_{x \to 0^+} \int_0^\pi \frac{t}{x}F\left(\frac{t}{x}\right) \cdot \frac{\sin t}{t}\mathrm{d}t \\
&\leqslant \varlimsup_{x \to 0^+} \int_0^\eta \frac{t}{x}F\left(\frac{t}{x}\right) \cdot \frac{\sin t}{t}\mathrm{d}t + \varlimsup_{x \to 0^+} \int_\eta^\pi \frac{t}{x}F\left(\frac{t}{x}\right) \cdot \frac{\sin t}{t}\mathrm{d}t \\
&\leqslant \sup_{y \in [0,+\infty)} yF(y) \cdot \int_0^\eta \frac{\sin t}{t}\mathrm{d}t + \varlimsup_{x \to 0^+} \sup_{y \in \left[ \frac{\eta}{x}, \frac{\pi}{x} \right]} yF(y) \cdot \int_0^\pi \frac{\sin t}{t}\mathrm{d}t \\
&= \sup_{y \in [0,+\infty)} yF(y) \cdot \int_0^\eta \frac{\sin t}{t}\mathrm{d}t,
\end{align*}
这里最后一个等号用到了\eqref{eq::23945789023589028226}. 由$\eta$任意性得
\[
\varlimsup_{x \to 0^+} \int_0^\infty F(t)\sin(tx)\mathrm{d}t \leqslant 0.
\]
结合\eqref{eq::23945789023589028234}即得\eqref{eq::23945789023589028227}.
\end{enumerate}
\end{proof}

\begin{example}
假设函数$g(x): (0,+\infty) \to (0,+\infty)$可导, 证明: 存在趋于正无穷的正数列$\{x_n\}$, 使得$g'(x_n) < g(2x_n)$, $n=1,2,\cdots$.
\end{example}
\begin{remark}
$\sup A=a<+\infty$已经包含了$A=\varnothing$的情况，因为
\begin{align*}
A=\varnothing \Longleftrightarrow \sup A=a=0.
\end{align*}
也可以用同样的方法先证明$A\ne \varnothing$。反证，若$A=\varnothing$，则
\begin{align*}
g'(x) \geqslant g(2x) >0,\forall x>0\Longrightarrow g\text{在}(0,+\infty)\text{上严格递增}.
\end{align*}
从而任取$x_0>1$，由Lagrange中值定理知，存在$\eta \in (x_0,2x_0)$，使得
\begin{align*}
g(2x_0)-g(x_0)=g'(\eta)x_0.
\end{align*}
进而
\begin{align*}
g(2\eta) \geqslant g(2x_0) > g(2x_0)-g(x_0)=g'(\eta)x_0 \Longrightarrow g'(\eta) < \frac{g(2\eta)}{x_0}<g(2\eta).
\end{align*}
故$\eta \in A$，这与$A=\varnothing$矛盾！
\end{remark}
\begin{proof}
记
\begin{align*}
A=\left\{ x>0:\,g'(x)-g(2x) <0 \right\},
\end{align*}
若$\sup A=a<+\infty$，则
\begin{align*}
g'(x) \geqslant g(2x) >0,\forall x>a\Longrightarrow g\text{在}(a,+\infty)\text{上严格递增}.
\end{align*}
对$\forall x>a$，由Lagrange中值定理知，存在$\xi \in (x,2x)$，使得
\begin{align*}
g(2x)-g(x)=g'(\xi)x.
\end{align*}
于是
\begin{align*}
g(2x)-g(x)=g'(\xi)x \geqslant g(2\xi)x \geqslant g(2x)x \Longleftrightarrow g(2x)(1-x) \geqslant g(x).
\end{align*}
令$x\rightarrow +\infty$，得到矛盾！故$\sup A=+\infty,$因此存在趋于正无穷的正数列$\{x_n\}$, 使得$g'(x_n) < g(2x_n)$, $n=1,2,\cdots$.
\end{proof}

\begin{example}
若$f(x)\in C^2[0,1]$, $f'(0)=0$, $|f''(x)|\leqslant |f(x)-f(0)|$. 证明：$f(x)$在$[0,1]$上为常值函数。
\end{example}
\begin{proof}
显然，存在$x_1\in(0,1]$，使得$f(x_1)\ne f(0)$. 否则$f(x)\equiv f(0)$，结论成立. 再由$f$的连续性知，可取$x_0\in(0,1]$满足
\begin{align*}
f(x_0)=\max\limits_{x\in[0,1]}|f(x)-f(0)|.
\end{align*}
由Taylor定理知，$\exists \xi\in(0,x_0)$，使得
\begin{align*}
f(x_0)=f(0)+f'(0)x_0+\frac{f''(\xi)}{2}x_0^2
=f(0)+\frac{f''(\xi)}{2}x_0^2.
\end{align*}
从而
\begin{align*}
|f''(\xi)|=\frac{2|f(x_0)-f(0)|}{x_0^2}
>|f(x_0)-f(0)|\geqslant |f(\xi)-f(0)|,
\end{align*}
这与题设矛盾! 故结论得证.
\end{proof}

\begin{example}
判断下述命题是否成立，并说明理由：若 $f$ 在 $\mathbb{R}$ 上可导且有界，并且
\begin{align*}
|f'(x)| < 1, \quad x \in \mathbb{R},
\end{align*}
则存在常数 $M < 1$，使得对任意 $x \in \mathbb{R}$，均有
\begin{align*}
|f(x) - f(0)| \leqslant M|x|.
\end{align*}
\end{example}
\begin{proof}
命题成立,证明如下:令
\begin{align*}
g(x) \triangleq 
\begin{cases}
\dfrac{f(x)-f(0)}{x}, & x \ne 0, \\
f'(0), & x = 0.
\end{cases}
\end{align*}
显然 $g \in C(\mathbb{R})$. 由 Lagrange 中值定理和题设知，对 $\forall x \ne 0$，存在 $\xi_x$ 在 $0$ 与 $x$ 之间，使得
\begin{align*}
|g(x)| = |f'(\xi_x)| < 1.
\end{align*}
故 $\sup\limits_{x \ne 0} |g(x)| \leqslant 1$. 又因为 $|g(0)| = |f'(0)| < 1$，所以只需证明
\begin{align*}
\sup\limits_{x \ne 0} |g(x)| < 1.
\end{align*}
反证，若 $\sup\limits_{x \ne 0} |g(x)| = 1$，则存在点列 $\{x_n\}_{n=1}^{\infty} \subset \mathbb{R} \setminus \{0\}$，使得
\begin{align}
\lim_{n \to \infty} |g(x_n)| = 1. \label{limg11}
\end{align}
由 $f$ 在 $\mathbb{R}$ 上有界知，$\exists K > 0$，使得
\begin{align}
|g(x_n)| = \frac{|f(x_n) - f(0)|}{|x_n|} \leqslant \frac{|f(x_n)| + |f(0)|}{|x_n|} \leqslant \frac{2K}{|x_n|}. \label{bound22}
\end{align}
若 $\varlimsup\limits_{n \to \infty} |x_n| = +\infty$，则由 \eqref{limg11}、\eqref{bound22} 式可得
\begin{align*}
1 = \varliminf_{n \to \infty} |g(x_n)| \leqslant \varliminf_{n \to \infty} \frac{2K}{|x_n|} = \frac{2K}{\varlimsup\limits_{n \to \infty} |x_n|} = 0,
\end{align*}
矛盾! 故 $\varlimsup\limits_{n \to \infty} |x_n| = a < +\infty$. 从而存在子列 $\{x_{n_k}\}_{k=1}^{\infty} \subset \{x_n\}_{n=1}^{\infty}$，使得 $\lim\limits_{k \to \infty} x_{n_k} = a$.
\begin{enumerate}[(i)]
\item 若 $a = 0$，则由 $|g|$ 的连续性可得
\begin{align*}
\lim_{k \to \infty} |g(x_{n_k})| = |g(0)| = |f'(0)| < 1,
\end{align*}
这与 \eqref{limg11} 式矛盾!

\item 若 $a \ne 0$，则由 $|g|$ 的连续性和 \eqref{limg11} 式可得
\begin{align}
1 = \lim_{k \to \infty} |g(x_{n_k})| = |g(a)| = \frac{|f(a) - f(0)|}{|a|}. \label{ga}
\end{align}
由 Lagrange 中值定理可知，存在 $\eta$ 在 $0$ 与 $a$ 之间，使得
\begin{align*}
\frac{|f(a) - f(0)|}{|a|} = f'(\eta) < 1,
\end{align*}
这与 \eqref{ga} 式矛盾!
\end{enumerate}
综上可知 $\sup\limits_{x \ne 0} |g(x)| < 1$.
\end{proof}

\begin{example}
设 $f$ 是 $\mathbb{R}$ 上有界或者有下界的函数，且满足对某个 $a>0$ 成立
\begin{align*}
f(x)+a\int_{x-1}^{x}f(y)\,\mathrm{d}y
\end{align*}
为常数，证明 $f$ 为常数。
\end{example}
\begin{proof}
不妨设 $f$ 有下界，否则用 $-f$ 代替 $f$. 再不妨设 $\inf_{x\in\mathbb{R}}f=0$，否则用 $f-\inf_{x\in\mathbb{R}}f$ 代替 $f$. 由题设可记
\begin{align*}
C\triangleq f(x)+a\int_{x-1}^{x}f(y)\,\mathrm{d}y,
\end{align*}
其中$C$ 为常数.从而
\begin{align*}
f(x)=C-a\int_{x-1}^{x}f(y)\,\mathrm{d}y.
\end{align*}
由于上式右端关于$x$可导，故 $f$ 可导,且
\begin{align*}
f'(x)=-a\bigl(f(x)-f(x-1)\bigr)\iff f'(x)+af(x)=af(x-1).
\end{align*}
由于 $\inf_{x\in\mathbb{R}}f=0$，故 $f(x)\geqslant 0$ 对一切 $x\in\mathbb{R}$ 成立，从而 $af(x-1)\geqslant 0$. 于是
\begin{align*}
e^{-ax}\bigl[e^{ax}f(x)\bigr]'=f'(x)+af(x)=af(x-1)\geqslant 0,
\end{align*}
因此 $e^{ax}f(x)$ 是单调递增函数。由此，对任意 $x\in\mathbb{R}$，有
\begin{align*}
C&=f(x)+a\int_{x-1}^{x}f(y)\,\mathrm{d}y =f(x)+a\int_{x-1}^{x}e^{ay}f(y)e^{-ay}\,\mathrm{d}y \\
&\leqslant f(x)+a\int_{x-1}^{x}e^{ax}f(x)e^{-ay}\,\mathrm{d}y =f(x)+a e^{ax}f(x)\cdot \frac{e^{-a(x-1)}-e^{-ax}}{a} \\
&=f(x)+f(x)(e^{a}-1)=f(x)e^{a}.
\end{align*}
所以 $f(x)\geqslant C e^{-a}$ 对一切 $x\in \mathbb{R}$成立。两边取下确界，得
\begin{align*}
0=\inf_{x\in\mathbb{R}} f(x)\geqslant C e^{-a}\Longrightarrow C\leqslant 0.
\end{align*}
另一方面，由于 $f(x)\geqslant 0$ 且 $a>0$，有
\begin{align*}
C=f(x)+a\int_{x-1}^{x}f(y)\,\mathrm{d}y\geqslant 0.
\end{align*}
故 $C=0$.于是
\begin{align*}
f(x)+a\int_{x-1}^{x}f(y)\,\mathrm{d}y=0.
\end{align*}
上式左边两项均非负，所以 $f(x)=0$ 对一切 $x\in \mathbb{R}$ 成立，即 $f\equiv 0$.
\end{proof}






\end{document}